% Section 6: Discussion (~0.75 pages)

\subsection{Disentangling Prompt Effects from Retrieval Effects}

Experiment~2 reveals that multi-step prevention, the apparent mechanism in Experiment~1, was a truncation artifact: only 1 of 1{,}200 episodes triggered multi-step reasoning.
The actual mechanism is \emph{output-driven efficiency}: the feedback rubric constrains LLM generation length directly (C4 output tokens 33\% below C1), acting as a metacognitive anchor that reduces verbose exploration rather than preventing reasoning spirals.

This has an immediate practical implication: any retrieval-augmented system can reduce token consumption by appending a structured feedback rubric, even without adaptive retrieval, a result independent of our Thompson Sampling contribution (see Appendix Figures~\ref{fig:token_trend} and~\ref{fig:theme_heatmap}).
The effect scales with task complexity (12\% for easy, 21\% for medium and hard) and operates consistently across all themes.

Experiment~2 also reveals a \emph{second} mechanism that was invisible in the small-library setting: retrieval differentiation.
With Jaccard overlap dropping from ${>}$0.90 to 0.52--0.59 (Table~\ref{tab:v3_results}), weight presets now produce genuinely different skill sets, and C3's retrieval quality gains (NDCG@5 $+$26.3\%, MRR $+$35.9\% vs.\ C1) break through the saturation ceiling diagnosed in Experiment~1.
The limitation was not fundamental to the R/I/V decomposition, but rather a consequence of insufficient library diversity.

\subsection{Input Assessment vs.\ Output Assessment}

Our design deliberately avoids output-level self-assessment~\citep{pan2024feedback}.
The empirical evidence (no detectable rating drift over 579 iterations in a prior deployment, 1{,}000 episodes in Experiment~1, and 1{,}200 episodes in Experiment~2) supports this choice.
Rating distributions in the first 50 episodes are statistically indistinguishable from the final 50 episodes across all dimensions ($p > 0.3$, Kolmogorov--Smirnov test; see Appendix Figure~\ref{fig:rating_stability}).
Beyond stability, the ratings are \emph{informative}: composite Likert scores correlate positively with ground-truth NDCG@5 ($r{=}0.33$, $p{<}0.001$; $r{=}0.44$ for C2 and C3), confirming that higher ratings correspond to genuinely better retrieval.
\citet{khalaf2025itrh} recently extended the ICRH finding to inference-time reward hacking, where best-of-$n$ sampling and inference-time alignment degrade quality; our input-assessment design is robust to both variants.

The principle is transferable: any system using LLM self-assessment in a feedback loop should evaluate upstream signals (inputs, context, retrieval) rather than downstream outputs to avoid self-referential reward inflation.

\subsection{When Does Weight Learning Matter?}

\begin{table}[t]
\centering
\small
\begin{tabular}{lrr}
\toprule
\textbf{Design / Metric} & \textbf{Exp.~1 (v2)} & \textbf{Exp.~2 (v3)} \\
\midrule
Skills              & 50 generic      & 205 domain-specific \\
Weight presets      & 5               & 12 \\
Tasks $\times$ conds & 250 $\times$ 4 & 300 $\times$ 4 \\
\texttt{max\_tokens} & 4{,}096        & 16{,}384 \\
\midrule
C1 mean tokens      & 6{,}481        & 5{,}517 \\
Multi-step rate (C1) & 11.6\%         & 0.3\% \\
\midrule
Jaccard C1--C2      & 0.904           & 0.585 \\
Jaccard C1--C3      & 0.909           & 0.517 \\
Jaccard C1--C4      & 0.848           & 0.573 \\
C3 NDCG@5           & ---             & 0.469 \\
C3 MRR              & ---             & 0.411 \\
\midrule
C2/C3 convergence   & importance      & pure\_relevance \\
C4 convergence      & none            & weak (importance) \\
Primary mechanism   & multi-step prev.\ & output compression \\
\bottomrule
\end{tabular}
\caption{Design changes and outcomes across experiments. Experiment~2's expanded library and de-saturated dimensions enable weight learning to produce differentiated retrieval (Jaccard $<$ 0.60), while eliminating truncation cascades reveals output compression as the primary efficiency mechanism.}
\label{tab:v2v3_comparison}
\end{table}

Experiment~1 shows that weight learning is inert with small, generic skill libraries where dimensions are saturated (Table~\ref{tab:v2v3_comparison}).
The mechanism is straightforward: Thompson Sampling's posterior width reflects uncertainty about each arm's quality, but when all presets retrieve nearly identical skill sets (Jaccard $> 0.90$), the observed reward variance is dominated by noise rather than genuine preset differences.
The posterior converges, but to a distinction without a difference.

Experiment~2 establishes the boundary condition.
With 205 domain-specific skills and 12 weight presets, C3 converges to \texttt{pure\_relevance} ($E[\theta]{=}$0.749) within ${\sim}$50 tasks (changepoint at tasks 33--48), and its explanation embeddings appear to produce a higher-fidelity reward signal than C2's Likert ratings.
The threshold at which weight learning begins to matter requires two conditions to be met simultaneously: sufficient library diversity (enough domain-specific skills that weight changes produce different rankings) and sufficient dimension variance (de-saturating recency and importance so that the weight space is not effectively one-dimensional).
Experiment~1 failed the first condition; Experiment~2 satisfies both.

The C4 result in the expanded setting is informative: C4 \emph{does} converge (to \texttt{pure\_importance}, $E[\theta]{=}$0.472), but with a compressed recency dimension (mean${=}$0.136).
This represents a shift from Experiment~1, where C4 showed flat posteriors.
The feedback rubric focuses the LLM's output regardless of what the bandit does with the weights.
C4's token advantage is an output-driven efficiency effect: the rubric constrains the LLM's generation. C3's advantage is a retrieval engineering effect: the bandit learns to retrieve better skills, aided by higher-fidelity reward signals from explanation embeddings.
These two mechanisms appear largely independent and can be combined.

These observations yield a two-condition diagnostic for practitioners considering adaptive retrieval weights.
First, compute the pairwise Jaccard similarity of top-$k$ retrievals across a representative sample of weight configurations: if Jaccard $> 0.80$, weight learning cannot differentiate presets and the investment is wasted (see Appendix Figure~\ref{fig:jaccard_vs_tokens} for the relationship between overlap and token consumption).
Second, check dimension variance across the skill library: if any dimension has variance $< 0.10$ (as recency and importance did in v2 at ${\sim}0.95$), the scoring function collapses to relevance-only ranking regardless of weights.
The Experiment~1 to Experiment~2 transition demonstrates this diagnostic in action: addressing both conditions dropped mean Jaccard from $>$0.90 to 0.52--0.59 and produced the retrieval quality improvements reported in Section~\ref{sec:experiment_v3}.

In practice, start with the rubric-only baseline (C4), which captures output-driven efficiency without any bandit machinery.
Upgrade to adaptive weight learning only if the Jaccard diagnostic shows sufficient skill diversity and the dimension variance check confirms a meaningful weight space.
Thompson Sampling itself is cheap (12 Beta posteriors, episode-wise updates), but the surrounding instrumentation (parsing ratings, computing diagnostics, validating convergence) is not, and should be justified by expected retrieval quality gains.
Where both diagnostic conditions are met, the investment pays off quickly: C3 converges in ${\sim}$50 tasks, and the two mechanisms (output efficiency + retrieval adaptation) can be combined.

\subsection{Transferability Beyond RAG}

While our experiments operate within vector-based RAG, the two primary findings are not structurally dependent on the retrieval architecture:

\begin{enumerate}
\item \textbf{Metalevel input assessment} (evaluating retrieval inputs rather than task outputs to avoid reward hacking) applies to any system that uses LLM self-assessment in a feedback loop.
The absence of rating drift across 2{,}200+ episodes is a property of the reward design, not the retrieval architecture.

\item \textbf{Output-driven efficiency} (structured feedback rubrics reducing output tokens by 33\%) operates through prompt structure, not retrieval.
Any system that appends a structured evaluation rubric may observe similar output compression.
\end{enumerate}

The retrieval weight optimization itself is specific to the Park et al.\ (2023) paradigm and is most valuable for large-scale deployments (thousands of documents) where the corpus exceeds context-window limits and fine-tuning lacks sufficient training data.
